\section{Materials and Methods}

\subsection{Device fabrication}
\label{ssec:dev_fab}
The memristive devices used in this study are fabricated on top of foundry-based MOS transistors, namely one-transistor~/~one-resistor (1T-1R), allowing the control of the compliance current $I_\mathrm{C}$~\cite{bricalli2018TED-I}. The \gls{be} consists of a 70\,nm\,$\times$\,70\,nm graphitic carbon pillar, that was already demonstrated to be a good electrode material for both volatile and non-volatile devices thanks to its stability and inert behavior \cite{covi2021TED-I}. The oxide layer and the \gls{te} are fabricated by e-beam evaporation. The oxide is a 10\,nm HfO$\mathrm{_2}$ active layer and the \gls{te} is a 100\,nm thick Ag layer. Both HfO$\mathrm{_2}$ and Ag are deposited at room temperature and without breaking the vacuum (pressure 3$\times$10$\mathrm{^{-6}}$\,mbar).


\subsection{Electrical setup for device characterization}

The electrical setup is the same as described in \cite{covi2021TED-I}, where the device is connected either to a semiconductor device parameter analyzer (DC characterization) or to a waveform generator and an oscilloscope (pulsed characterization). DC characterization is carried out using a HP~4156C semiconductor device parameter analyzer. A sweep from 0\,V to 1.5\,V is applied to the \gls{te} of the device while the \gls{be} is grounded. The characterization includes DC sweeps with different compliance current (I$\mathrm{_C}$), as in Fig.~\ref{fig:Supfig_01}a. The pulsed characterization as well as the working memory experiment are carried out using a TTI~TGA12104 arbitrary waveform generator. The voltage waveforms are applied to the \gls{te}. To measure the current, a LeCroy Waverunner 640Zi oscilloscope is connected at the \gls{be} side, and the voltage drop across a 50\,$\Omega$ series resistance is probed. MATLAB\textsuperscript{\textregistered} software is used for the data analysis and the control of the instruments.

The temporal dynamics are studied by first applying a semi-triangular pulse with 10\,ms pulse duration and 5\,V amplitude to induce the filament formation and then monitoring the state of the filament with constant -150\,mV bias. To tune the temporal behavior, the I$\mathrm{_C}$ is changed from 10\,\textmu A to 70\,\textmu A. For each I$\mathrm{_C}$, 100 experiments are carried out. To avoid possible interference due to the previous cycles, the value of I$\mathrm{_C}$ is changed in random order.

The switching probability (Fig.\ref{fig:fig2}D) is studied by applying 100 pulses for each combination of voltage amplitude - pulse duration. The conditions are applied in random order. The impact of the number of pulses (Fig.\ref{fig:fig2}D and E) is analyzed with the same methodology, selecting the order of the combination of voltage amplitude - number of pulses randomly. Each combination is applied 100 times.

For the \gls{wm} store and recall experiment, each probability-frequency-pattern condition is repeated 10 times. The voltage bias is applied only at the end of the experiments to check the retention capabilities.

\subsection{Fitting of device features}

Thanks to the stochasticity of the switching mechanism, the probability P$_{ON}$ for a single pulse of given amplitude (Fig.~\ref{fig:fig2}D) follows a normal distribution and thus is fit using a cumulative distribution function:

\begin{equation} 
\label{eqPON}
{\Large\displaystyle{P_{ON}(V) = \frac{1}{2}\left(1+erf\left(\frac{V-\mu}{\sigma\sqrt{2}}\right)\right)}}
\end{equation}

where the average value $\mu$ and variance $\sigma$ are fitting parameters of the device and the pulse duration, respectively. Table~\ref{tab:tab1} collects the data of the device presented in Fig.~\ref{fig:fig2}D. Each point in the figure corresponds to 100 measurements. %\erika{@Saverio, can you please specify if it's a single device measurement and how many measurements you did to collect these data? ---> [saverio] if we want to show the device-to-device variability, I have an image which compares the 5 used devices. It was also in the first version of the supplementary material}


\begin{table}[ht]
    \centering
    \caption{Values of $\mu$ and $\sigma$ calculated for different pulse time widths.}
    \begin{tabular}{|c|c|c|}
        \hline
        \textbf{Time width [ms]} & \textbf{Average $\mu$ [V]} & \textbf{Variance $\sigma$ [V]} \\
        \hline
        0.05 & 2.31 & 0.38 \\
        \hline
        0.10 & 2.11 & 0.33 \\
        \hline
        0.15 & 1.86 & 0.30 \\
        \hline
        0.50 & 1.73 & 0.22 \\
        \hline
        1.00 & 1.21 & 0.16 \\
        \hline
        2.00 & 0.61 & 0.15 \\
        \hline
        5.00 & 0.59 & 0.11 \\
        \hline
    \end{tabular}
    \label{tab:tab1}
\end{table}

The switching probability, for a given amplitude, as a function of the number of pulses is fit with the mathematical model: 

\begin{equation} 
\label{eqPON_N}
P_{ON}(N) = 1 - (1 - P_{ON}(1))^N\quad \mathrm{where} \quad P_{ON}(1) = P_{ON}(V)
\end{equation}

In the \gls{wm} store and recall experiments, unless otherwise noted, the following conditions apply: no read voltage between spikes during the store and the recall phases. A read voltage of -150\,mV is applied at the end of the experiment. The current compliance is I$\mathrm{_C}$ = 17\,\textmu A, which gives an average retention time of 28\,ms.

\subsection{STP synapse model}
\label{sec:stpwmm}

To perform the computer simulations we developed a simplified phenomenological model that qualitatively reproduces the the behavior of the memristive devices. The model captures the switching probabilities and variable retention time of the devices. Each synapse $i$ was modeled with a binary internal state variable $x_i$ that denotes either the low ($x_i=1$) or the high resistance state ($x_i=0)$. The device resistance was $r_1$ and $r_0$ in the low and high resistance state, respectively. To model switching probabilities we assigned a parameter $\rho_i \in [0,1]$ to every synapse. Upon arrival of a pre-synaptic input spike, $x_i$ was set to 1 with probability $\rho_i$.

To model the trial-by-trail variability of the retention times we adopted a Lognormal distribution for the retention times $t_{ret}$, namely:
\begin{equation}
  \begin{split}
    t_{ret} \; &\sim\; \mathrm{Lognormal}( t_{ret} \;|\; \mu_i, \sigma_i )
    =\; \frac{1}{\sqrt{2}\pi\sigma_i \, t_{ret}} \, \exp \left( -\frac{(\log(t_{ret}) - \mu_i)^2}{2\sigma_i^2} \right) \;.
  \end{split}
  \label{eqn:lognormal}
\end{equation}
The parameters $\mu_i$ and $\sigma_i$ in Eq.~\ref{eqn:lognormal} were adjusted to fit the device properties. Supplementary figure~\ref{fig:Supfig_lognormal_fit} shows example model fits to experimental data for different compliance currents. In the simulations we used $\mu =  7.24$ and $\sigma_i = 0.82$, if not stated otherwise, which corresponds to a mean retention time of $\sim1.5$\,s. Whenever a synapse was set to the low resistance state ($x_i=1$) the retention time $t_{ret}$ was  drawn from Eq.~\ref{eqn:lognormal}. After the simulation time $t$ exceeded $t+t_{ret}$ the synapse spontaneously returned to the high-resistance state.

\subsection{Working memory network model}
\label{sec:network}

To reproduce the working memory model introduced in \cite{mongillo2008synaptic}, we used a recurrent network of 8000 excitatory and 2000 inhibitory neurons. Connection probabilities between these neurons were as in \cite{mongillo2008synaptic}. Supplementary Figure~\ref{fig:Supfig_08} illustrates the detailed network structure, connection probabilities and baseline synaptic conductances. Each memory item was represented by a population of 800 excitatory neurons, which were randomly chosen from the recurrent network for each of the 5 memory items. Synaptic weights between these neurons were 5 times stronger (0.5\,mV) than between other neurons (0.1\,mV). All inhibitory synapses were static (no STP) and had a strength of -0.2\,mV. To store and recall the working memory, 1000 input neurons were set up for each memory item. Input neurons were only connected to one of the memory item populations in the network. Input neurons fired at low baseline firing rates of 0.1\,Hz. To store a memory item, firing rates of input neurons corresponding to one memory item were elevated 10-fold. During recall all input neurons were set to unspecific elevated activation.

All neurons used the leaky integrate and fire (LIF) model with biologically plausible parameters \cite{gerstner2014neuronal}. The LIF neuron is a spiking point neuron model that transiently integrates synaptic inputs using a leaky membrane potential. The membrane potential $u(t)$ follows the dynamics
\begin{equation}
\frac{d\,u}{d\,t}  \;=\; - \frac{1}{\tau_m} \left(u(t) - u_0 \right) \;+\; i(t)\;,
\label{eq:membrane-potential}
\end{equation}
where $u_0$ is the membrane resting potential $\tau_m$ is a time constant. $i(t)$ is the input into the neuron denoting the summed effect of afferent synapses. If the membrane potential crosses a firing threshold $\vartheta$ at time $t^f$ a spike is emitted and the membrane potential is reset immediately after
\begin{equation}
 u(t^f+\Delta t) = u_r\;.
 \label{eq:reset_behavior}
\end{equation}
After a spike the neuron is inactive for a brief refractory time. The firing threshold was set to 20\,mV and refractory time to 2\,ms. Membrane time constants $\tau_m$ were 15\,ms and 10\,ms, resting potentials $u_0$ 16~mV and 13\,mV, for excitatory and inhibitory neurons, respectively. Independent unit variance Gaussian noise with mean 0.5775\,mV were injected into each excitatory, and with mean 0.5275\,mV into inhibitory neurons. As in \cite{mongillo2008synaptic} the mean of the Gaussian noise was precisely tuned to enable multistability in the network. Example network dynamics are shown in Fig.~\ref{fig:tsodyks-wm}D, for a store-recall experiment over 8 seconds.

\subsection{Associative symbolic working memory model}
\label{sec:associative-memory}

For the associative working memory model in Fig.~\ref{fig:symbolic-working-memory} we developed a network architecture that was capable to store bidirectional short-term associations between items of different categories. The associative memory model consisted of 8 memory neurons that were connected through bidirectional associative connections as illustrated in Fig.~\ref{fig:symbolic-working-memory}A. The memory neurons encoded for the three memory categories, shape, texture and color. Each memory neuron exclusively received input from one input neuron to trigger store and recall events. Memory neurons were implemented using the LIF neuron model, with firing threshold of 1.4\,V and membrane time constant of 15\,ms. Each group of neurons corresponding to one memory category was augmented with a single inhibitory neuron that provided lateral feedback to enforce mutually exclusive activation of one memory item at a time. Inhibitory neurons were LIF neurons with threshold of 1.4\,V and membrane time constant of 10\,ms. All excitatory connections had a strength of 1.5\,V and inhibitory feedback was -1.5\,V.

Model memristive devices were used inside the bidirectional connections to store the association between a specific pair of memory neurons (which we call here the specific neurons). Exactly one model device was used per association. Memristive devices were augmented with two auxiliary threshold gates to route the input and outputs during store and recall cycles. Gates were put in series with the devices, one before (input gate) and one after (output gate). Input gates received excitatory input from the two specific memory neurons the association corresponded to. In addition input gates received inhibitory input from all other memory neurons. Output gates recursively connected back to specific memory neurons through excitatory connections. Storage of memory items was not dependent on the activity of the neurons and threshold gates (see Fig.~\ref{fig:symbolic-working-memory}), but retained solely in the hidden state of the memristive devices. Model parameters of memristive devices were fit to experimental data ($\mu$=7.24, $\sigma$=0.82, corresponding to ca. 1.4\,s mean retention time). Switching probabilities $\rho$ were 0.2. The synaptic conductance for the low resistance state was chosen to elicit a voltage pulse of $1.0\,V$ in the post-synaptic neuron. The high resistance $r_0$ was set to be 20 times higher than the low resistance $r_1$.

\subsection{Details to software simulations}

Simulations of the working memory model with volatile memristive devices were done in python (3.8) using a custom implementation of the memristive synapse model that was developed for the NEST~2.14 simulation environment \cite{Peyser:838729}. The simulation time step was 1\,ms. Neuron dynamics were simulated using the current-based leaky integrate-and-fire model is available in NEST. A custom model of the STP model outlined in Section~\ref{sec:stpwmm} was implemented based on existing synapse models. Data analysis used the numpy and matplotlib python packages in version 1.23.2 and 3.5.3, respectively. The code will be made available upon publication.
